\documentclass[11pt,a4paper]{article}
\usepackage[T1]{fontenc}
\usepackage{mathpazo}
\usepackage[margin=26mm,headheight=14pt]{geometry}
\usepackage{amsmath,amssymb,amsthm,mathtools}
\usepackage{microtype,booktabs,tabularx,array,enumitem}
\usepackage{tocloft}
\setlength{\cftbeforesecskip}{3pt}
\usepackage[dvipsnames]{xcolor}
\usepackage{fancyhdr}
\usepackage{needspace}
\definecolor{ink}{HTML}{194E55}
\usepackage[colorlinks=true,linkcolor=ink,citecolor=ink,urlcolor=ink,
  pdftitle={Spectral Transfer for Finite-Field Constraints and Compressed Reflection-Positive Laws},
  pdfauthor={OpenAI-assisted research report prepared for Vladimir Reshetnikov}]{hyperref}
\usepackage{bookmark}
\usepackage{xurl}
\pagestyle{fancy}
\fancyhf{}
\fancyhead[L]{\small Spectral transfer and reflection-positive laws}
\fancyhead[R]{\small 7 October 2026}
\fancyfoot[C]{\thepage}
\setlength{\emergencystretch}{2em}
\setlist{itemsep=3pt,topsep=5pt}
\numberwithin{equation}{section}
\newtheorem{theorem}{Theorem}[section]
\newtheorem{lemma}[theorem]{Lemma}
\newtheorem{proposition}[theorem]{Proposition}
\newtheorem{corollary}[theorem]{Corollary}
\theoremstyle{definition}
\newtheorem{definition}[theorem]{Definition}
\newtheorem{example}[theorem]{Example}
\newtheorem{question}[theorem]{Question}
\theoremstyle{remark}
\newtheorem{remark}[theorem]{Remark}
\newcommand{\E}{\mathbb E}
\newcommand{\F}{\mathbb F}
\newcommand{\C}{\mathbb C}
\newcommand{\R}{\mathbb R}
\newcommand{\one}{\mathbf 1}
\newcommand{\op}{\mathrm{op}}
\newcommand{\Unif}{\operatorname{Unif}}
\newcommand{\Prob}{\mathbb P}
\newcommand{\supp}{\operatorname{supp}}
\newcommand{\diag}{\operatorname{diag}}
\newcommand{\spec}{\operatorname{spec}}
\newcommand{\norm}[1]{\left\lVert #1\right\rVert}
\newcommand{\abs}[1]{\left\lvert #1\right\rvert}
\newcommand{\sourcecommit}{\texttt{adc7f1241b42e322a6451854ab7e4b4c146bf78a}}
\title{\color{ink}\textbf{Spectral Transfer for Finite-Field Constraints}\\[4pt]
\Large and Compressed Reflection-Positive Laws}
\author{Research report prepared for Vladimir Reshetnikov\\
\small OpenAI-assisted mathematical development for the ProveIt collection}
\date{7 October 2026}
\begin{document}
\maketitle
\begin{abstract}
We prove that a single-edge spectral estimate passes directly to a product
of constraints on any simple bipartite graph, with no loss in the exponent
of the error. For edge kernels bounded by one, with constant models
$d_e\in[0,1)$ and spectral errors $\epsilon_e$, the resulting operator error is
at most $\sum_j\epsilon_{\pi(j)}\prod_{i<j}d_{\pi(i)}$ for every edge
ordering $\pi$. The best such ordering is obtained by sorting
$\epsilon_e/(1-d_e)$. In the homogeneous case this is
$\epsilon(1-d^m)/(1-d)$, uniformly in the number of vertices and edges.
For directed nonzero-square constraints over $\F_p$, this replaces the
$p^{-1/8}$ list estimate in the September 2026 OpenAI square-difference
manuscript by an explicit $O(p^{-1/2})$ estimate. The exponent is optimal
uniformly over all patterns.

We then reprove and optimize the finite reflection-positive construction
used by that manuscript. Its 161 permutation blocks can be replaced by
17, reducing the number of labels from $(24!)^{161}$ to $(24!)^{17}$,
while preserving the designated square-difference cycle and an exceptional
mass at most $p^{-4}$. For the optimized construction, the exceptional
mass has exact order $p^{-9/2}$ and the conditional prediction operator
has exact order $p^{-1/4}$. A collision argument proves the matching lower
bound. The construction extends to every fixed cycle length at least
three and every fixed-index multiplicative subgroup of a finite field.
These are finite, local theorems; no improved global exponent for
square-difference-free subsets of the integers is asserted.
\end{abstract}
\vspace{2mm}
\noindent\textbf{Keywords.} Spectral counting; finite fields; quadratic residues;
reflection positivity; conditional expectation; quantitative additive combinatorics.
\setcounter{tocdepth}{1}
\begingroup\small
\tableofcontents
\endgroup
\newpage

\section{Research setting, selection, and scope}
\label{sec:setting}

\subsection{The repository intersection}

The October 2026 release of \texttt{openai/math} contains a large collection
of research manuscripts, together with formal proof material for selected
statements. Its own introduction distinguishes the manuscripts from their
individual verification status \cite{catalogue}. Two closely related
families are especially relevant here: Family 159, on quasipolynomial
Szemer\'edi bounds, and Family 182, on polynomial-difference avoidance.
The latter includes \emph{A power saving for square-difference-free sets}
\cite{square}. Its finite-field construction is the source of the present
problem. We use the source at the immutable commit
\begin{center}\small\sourcecommit.\end{center}

The ProveIt snapshot reviewed here is
\texttt{fe7165a3cda1f9182b2a258cf13431a906c4b0c5}.
Recent work in its Gowers--Szemer\'edi project carefully separates local
inequalities, sharp constants, counterexamples to intermediate interfaces,
and complete global conclusions. The intake records for Reports 288--299
include hereditary additive-energy rigidity, sharp weighted extremizers,
fourth-norm estimates, and failures of an overly universal lifting
statement \cite{proveit,energyintake,refinementsintake}. This is a useful
methodological intersection: identify a precise finite estimate, improve
it without changing its quantifiers, and trace the improvement only through
arguments whose hypotheses have been checked.

This is a selective review of that intersection, not an audit of all the
manuscripts in either repository. In particular, the global theorems
reported in Families 159 and 182 are not needed as hypotheses of any
theorem below. The finite square-constraint and tuple constructions are
read and reproved directly.

\subsection{The specific loss being removed}

For a simple directed bipartite pattern with $m$ edges, the list-mixing
lemma in the source's finite-field graph section \cite{square}
(label \texttt{lem:list-mixing}) bounds the normalized operator error by
\begin{equation}
 (32m)^{1/4}p^{-1/8}.
 \label{eq:old-bound}
\end{equation}
It first bounds a four-cycle moment using scalar graph counting and then
takes a fourth root. The single-edge estimate already has order
$p^{-1/2}$. For a bipartite list kernel, whole-list $L^2$ test functions
can be retained while one edge is exposed. This makes the fourth-root
step unnecessary.

The distinction between whole-list tests and products of one-coordinate
tests matters. An operator bound concerns arbitrary $F(x_1,\ldots,x_a)$
and $G(y_1,\ldots,y_b)$ in $L^2$. Theorem~\ref{thm:replacement} treats
these arbitrary functions. It is therefore a direct replacement for the
operator lemma in the source, rather than merely a scalar counting result.

\subsection{Contributions and attribution}

The elementary ingredients---telescoping products, the spectral form of
mixing, Gauss sums, and positivity of Gram matrices---are classical.
For background on graph counting and additive combinatorics see
\cite{zhao}. The edge-exposure arguments of Conlon, Fox, and Zhao
\cite[Section 3]{cfz} are particularly close methodological precedents.
Examples of the surrounding quantitative progression and
square-difference literature are \cite{lss,bloommaynard}. We do not claim
that telescoping or spectral counting is new. The concrete contributions
established here are the following.

\begin{enumerate}
\item A general whole-list spectral replacement statement with geometric
      coefficients and its exactly optimal replacement order.
\item The explicit $p^{-1/2}$ replacement of \eqref{eq:old-bound}, including
      the exact single-edge constant for both odd prime congruence classes.
\item A reparameterized, fully proved finite tuple construction with 17
      blocks, with a general tradeoff between marking depth and exceptional
      mass, also valid for fixed-index multiplicative subgroups.
\item A matching collision lower bound showing that the optimized
      construction has conditional norm $\Theta(p^{-1/4})$.
\end{enumerate}

The permutation-block architecture and its marked-cut positive matrix
come from \cite[\texttt{build/sections/tuples.tex}]{square}; they are
explicitly credited and reproved in Sections~\ref{sec:construction}--\ref{sec:mixing}.
The reduction in depth and the sharp conditional lower bound are the new
developments of that architecture. Priority beyond the sources searched
has not been exhaustively established. All the stated mathematical
claims have proofs here; none is represented as a new Lean formalization.

\section{A whole-list spectral replacement theorem}
\label{sec:replacement}

Every finite probability space below has its specified probability
measure; products always mean independent product measures. Zero-mass
points can be removed. For a complex kernel $B:X\times Y\to\C$, set
\[
 T_Bg(x)=\E_y B(x,y)g(y),\qquad
 \norm{B}_{\op}=\norm{T_B}_{L^2(Y)\to L^2(X)}.
\]
Equivalently,
\begin{equation}
 \abs{\E_{x,y}\overline{f(x)}B(x,y)g(y)}
 \le \norm{B}_{\op}\norm f_2\norm g_2.
 \label{eq:duality}
\end{equation}
A constant kernel $c$ is the rank-one operator $g\mapsto c\E g$.
For uniform spaces of sizes $a,b$, the matrix representing $T_B$ in
orthonormal point-mass bases is $B/\sqrt{ab}$, not the unnormalized
matrix $B$.

Let $I,J$ be disjoint finite index sets. Give each $i\in I$ a probability
space $X_i$ and each $j\in J$ a probability space $Y_j$. Write
$X=\prod_iX_i$ and $Y=\prod_jY_j$. A pattern $H$ is a subset of
$I\times J$; thus there is at most one edge for each pair of coordinates.
Assign a kernel $W_e:X_i\times Y_j\to\C$ to every $e=(i,j)$ and define
\[
 K_W(x,y)=\prod_{e=(i,j)\in H}W_e(x_i,y_j).
\]

\begin{lemma}[One exposed edge]
\label{lem:exposed}
Fix $e=(i,j)\in H$. On every other edge choose a kernel of absolute
value at most one, and put an arbitrary kernel $B$ on $e$. Then the
resulting whole-list kernel $L$ satisfies
\[
 \norm L_{\op}\le\norm B_{\op}.
\]
\end{lemma}
\begin{proof}
Test $L$ against arbitrary $F\in L^2(X)$ and $G\in L^2(Y)$.
Fix $u=x_{I\setminus\{i\}}$ and $v=y_{J\setminus\{j\}}$.
All factors other than $B(x_i,y_j)$ separate into three groups: factors
depending on $x_i$, factors depending on $y_j$, and factors depending
only on $u,v$. No other factor depends on both $x_i$ and $y_j$, by
simplicity of the pattern. The conditional pairing has the form
\[
 C(u,v)\E_{x_i,y_j}\overline{F_u(x_i)}B(x_i,y_j)G_v(y_j),
\]
where the notation absorbs the incident edge factors into the slices,
and
\[
 |C(u,v)|\le1,\quad
 \norm{F_u}_2\le\norm{F(u,\cdot)}_2,\quad
 \norm{G_v}_2\le\norm{G(v,\cdot)}_2.
\]
The absorbed slices can depend on both $u$ and $v$; their displayed
upper bounds depend only on $u$ and only on $v$, respectively.
Applying \eqref{eq:duality}, then averaging the independent $u,v$, gives
\begin{align*}
 |\langle F,T_LG\rangle|
 &\le\norm B_{\op}
   \bigl(\E_u\norm{F(u,\cdot)}_2\bigr)
   \bigl(\E_v\norm{G(v,\cdot)}_2\bigr)\\
 &\le\norm B_{\op}\norm F_2\norm G_2.
\end{align*}
The last step is Cauchy--Schwarz on each remaining product space.
\end{proof}

\begin{theorem}[Spectral replacement]
\label{thm:replacement}
Assume $|W_e|\le1$, choose $d_e\in[0,1]$, and suppose
$\norm{W_e-d_e}_{\op}\le\epsilon_e$. If $m=|H|$ and $\pi$ is any
ordering of its edges, then
\begin{equation}
 \boxed{\quad
 \norm{K_W-\prod_e d_e}_{\op}
 \le S(\pi):=\sum_{j=1}^m\epsilon_{\pi(j)}
                    \prod_{i<j}d_{\pi(i)}.\quad}
 \label{eq:replacement}
\end{equation}
The assertion includes the empty pattern, with error zero. It has no
dependence on the dimensions of the coordinate spaces or on vertex
degrees.
\end{theorem}
\begin{proof}
Order the edges as $e_1,\ldots,e_m$. The exact telescoping identity is
\[
 \prod_{j=1}^mW_{e_j}-\prod_{j=1}^md_{e_j}
 =\sum_{j=1}^m\left(\prod_{i<j}d_{e_i}\right)
    (W_{e_j}-d_{e_j})\left(\prod_{i>j}W_{e_i}\right).
\]
Apply Lemma~\ref{lem:exposed} to each summand and use the triangle
inequality for the operator norm.
\end{proof}

\begin{corollary}[Homogeneous parameters]
\label{cor:homogeneous}
If all $d_e=d<1$ and all $\epsilon_e\le\epsilon$, then
\begin{equation}
 \norm{K_W-d^m}_{\op}
 \le\epsilon\frac{1-d^m}{1-d}.
 \label{eq:homogeneous}
\end{equation}
For $d=1$, the bound is $m\epsilon$.
\end{corollary}

\begin{corollary}[Comparison of two nonconstant models]
\label{cor:nonconstant}
If both $|W_e|\le1$ and $|V_e|\le1$, then
\[
 \norm{K_W-K_V}_{\op}
 \le\sum_{e\in H}\norm{W_e-V_e}_{\op}.
\]
\end{corollary}
\begin{proof}
Use the same telescoping identity with $V_e$ replacing $d_e$. All
factors other than the exposed difference are still bounded by one.
\end{proof}

\subsection{The exactly optimal replacement order}

\begin{proposition}[Ordering rule]
\label{prop:order}
Among all bounds $S(\pi)$ in \eqref{eq:replacement}, a minimum is
obtained by arranging the edges in nondecreasing order of
\[
 \frac{\epsilon_e}{1-d_e}\qquad(d_e<1).
\]
Edges with $d_e=1$ and $\epsilon_e>0$ may be put last. Edges with
$d_e=1$ and $\epsilon_e=0$ may be put anywhere. Ties may be ordered
arbitrarily. This optimizes this telescoping bound, not the actual
operator norm.
\end{proposition}
\begin{proof}
Consider adjacent edges $e,f$. Their contribution, after removing a
common nonnegative prefix product, is
$\epsilon_e+d_e\epsilon_f$. Reversing them gives
$\epsilon_f+d_f\epsilon_e$. Their difference is
\[
 \epsilon_e(1-d_f)-\epsilon_f(1-d_e).
\]
All subsequent terms have the same prefix product $d_ed_f$ in either
order. An adjacent inversion of the claimed ordering therefore never
increases the total. Sorting by successive interchanges proves the result,
including the stated $d=1$ cases. A zero prefix can produce additional
minimizers; no uniqueness assertion is needed.
\end{proof}

\subsection{Filters, scalar counts, and indispensable hypotheses}

If $a:X\to\C$ and $b:Y\to\C$ satisfy $|a|,|b|\le1$, multiplication
on the two sides is a contraction. Consequently
\[
 \norm{a(x)\bigl(K_W(x,y)-\prod_e d_e\bigr)b(y)}_{\op}
 \le S(\pi).
\]
This permits arbitrary bounded constraints entirely within either list.
The corresponding model is $a(x)(\prod_e d_e)b(y)$, not a constant
kernel. If one instead conditions on events of probabilities $\alpha$
and $\beta$ on the two sides, the operator bound under the conditional
product measures is at most $S(\pi)/\sqrt{\alpha\beta}$, by extending
the two tests by zero and keeping the normalization factors.

For an arbitrary simple directed graph on independent coordinate
variables, the scalar analogue of \eqref{eq:replacement} also holds.
Expose the two endpoints of one edge and apply \eqref{eq:duality} to
bounded one-variable functions. The remaining tests have norm at most
one. In particular,
\begin{equation}
 \abs{\E\prod_eW_e-\prod_ed_e}\le S(\pi).
 \label{eq:scalar}
\end{equation}
The whole-list operator assertion, in contrast, uses the specified
bipartition.

The hypotheses cannot be omitted indiscriminately. If the same
$0$--$1$ kernel $W$ occurs twice on one pair of coordinates, its product
is $W$, not a two-edge pseudorandom product. For the square kernel below,
$\norm{W-1/2}_{\op}\to0$, whereas
$\norm{W^2-1/4}_{\op}\ge|\E W-1/4|\to1/4$.
Thus repeated edges would invalidate a vanishing two-edge conclusion.
Likewise, product measure is essential. For
$B(x,y)=1_{x=y}$ on a uniform $n$-point space,
$\norm B_{\op}=1/n$, but $\E B(X,Y)=1$ if one imposes $Y=X$.
The independence used in the proof is a mathematical assumption, not
just a convenient notation.

\section{Finite-field residue kernels}
\label{sec:fields}

\subsection{Multiplicative subgroups over arbitrary finite fields}

Let $q$ be a prime power and let $H\le\F_q^\times$ have index
$s\ge2$. Set
\[
 A_H(t)=1_H(t),\qquad d_s=1/s,\qquad
 \delta_{q,s}=\frac{(s-1)\sqrt q+1}{sq}.
\]
The symbol $s$ here denotes the subgroup index; marking depth below is
denoted by $t$, and the size of a mark set by $j$.

\begin{lemma}[Single-edge subgroup estimate]
\label{lem:subgroup}
On the uniform additive group of $\F_q$,
\[
 \norm{A_H(y-x)-1/s}_{\op}\le\delta_{q,s}.
\]
The same bound holds for the reversed difference.
\end{lemma}
\begin{proof}
Choose a multiplicative character $\chi$ of exact order $s$ and
extend all multiplicative characters by zero at zero. For $t\ne0$,
\[
 A_H(t)=\frac1s\sum_{j=0}^{s-1}\chi(t)^j.
\]
Additive characters diagonalize convolution on $\F_q$. The eigenvalue
at the trivial additive character after subtracting $1/s$ is $-1/(sq)$.
For a nontrivial additive character $\psi$, the unnormalized Fourier
coefficient of $A_H$ is
\[
 \frac1s\left(-1+\sum_{j=1}^{s-1}
            \sum_{u\ne0}\chi(u)^j\psi(u)\right).
\]
Each nontrivial Gauss sum in this expression has absolute value $\sqrt q$.
Here is the required calculation. For nontrivial multiplicative $\xi$,
put $G(\xi,\psi)=\sum_{u\ne0}\xi(u)\psi(u)$. The substitution $u=av$
gives
\begin{align*}
 |G(\xi,\psi)|^2
 &=\sum_{a\ne0}\xi(a)\sum_{v\ne0}\psi((a-1)v)\\
 &=(q-1)-\sum_{a\ne0,1}\xi(a)=q,
\end{align*}
because $\sum_{a\ne0}\xi(a)=0$. Dividing the Fourier coefficient by
$q$ proves the estimate. Transposition preserves the operator norm.
\end{proof}

\begin{theorem}[Residue constraints between lists]
\label{thm:residuelists}
For any simple directed bipartite pattern with $m$ edges, let $K_H$
be the product of the indicators that the directed differences lie in
$H$. Directions can differ on different edges. Then
\begin{equation}
 \boxed{\quad
 \norm{K_H-s^{-m}}_{\op}
 \le(1-s^{-m})\left(q^{-1/2}+\frac1{(s-1)q}\right).
 \quad}
 \label{eq:subgrouplists}
\end{equation}
The same right-hand side bounds the scalar counting error for any
simple directed graph with no loops and no opposite edges.
\end{theorem}
\begin{proof}
Apply Corollary~\ref{cor:homogeneous} and \eqref{eq:scalar} with
$d=1/s$ and $\epsilon=\delta_{q,s}$.
\end{proof}

For fixed $s$ and a fixed graph, \eqref{eq:scalar} shows that its
constraint probability tends to $s^{-m}>0$. This elementary positivity
fact, with an explicit error, is all that is needed to normalize the
finite laws below. The bound itself is uniform in the pattern, even
when its size varies with $q$. It is an absolute-error estimate; a
relative-error conclusion requires $s^{-m}$ to dominate the displayed
error.

\subsection{Exact constants for nonzero squares}

For an odd prime $p$ write $Q_p^*$ for the nonzero squares and
$A_p=1_{Q_p^*}$. Define
\begin{equation}
 \Delta_p=
 \begin{cases}
  (\sqrt p+1)/(2p),&p\equiv1\pmod4,\\
  \sqrt{p+1}/(2p),&p\equiv3\pmod4.
 \end{cases}
 \label{eq:exactdelta}
\end{equation}

\begin{proposition}[Exact square-kernel norm]
\label{prop:exactsquare}
For every odd prime,
\[
 \norm{A_p(y-x)-1/2}_{\op}=\Delta_p.
\]
Consequently every simple directed bipartite square pattern satisfies
\begin{equation}
 \norm{K-2^{-m}}_{\op}
 \le2(1-2^{-m})\Delta_p
 \le(1-2^{-m})\frac{\sqrt p+1}{p}.
 \label{eq:newsquare}
\end{equation}
No uniform $o(p^{-1/2})$ estimate is possible for all nonempty patterns.
\end{proposition}
\begin{proof}
Let $\chi$ be the quadratic character and
$G=\sum_{u\ne0}\chi(u)e^{2\pi iu/p}$. The trivial Fourier eigenvalue
of the centered kernel is $-1/(2p)$. All its other eigenvalues are
\[
 \frac{\chi(k)G-1}{2p},\qquad k\ne0.
\]
The preceding Gauss-sum calculation gives $|G|=\sqrt p$, and complex
conjugation gives $\overline G=\chi(-1)G$. Thus $G$ is real for
$p\equiv1\pmod4$ and purely imaginary for $p\equiv3\pmod4$.
Both values of $\chi(k)$ occur. Taking the largest eigenvalue modulus
proves \eqref{eq:exactdelta}, without needing the sign of the real
Gauss sum. Formula~\eqref{eq:newsquare} follows from
Corollary~\ref{cor:homogeneous}. A pattern with a single edge has
norm exactly $\Delta_p\sim\tfrac12p^{-1/2}$, proving the final claim.
\end{proof}

The exact density is $(p-1)/(2p)$. One may center there instead of at
$1/2$, removing the trivial eigenvalue; the nontrivial eigenvalues and
their maximum are unchanged. Centering at $1/2$ is convenient because
it matches the source construction. The distinction is kept explicit
in all formulas.

\section{A general marking principle}
\label{sec:marking}

The next finite lemma explains why a better spectral exponent reduces
the number of blocks. It also separates the measure construction from
its bookkeeping.

Let $\Lambda$ be a set of $\ell=2t+1$ blocks, where $t\ge1$.
On a finite label space, suppose there are nonnegative measures
$\nu_S$, one for every $S\subseteq\Lambda$ with $|S|\le t+1$.
For each block there may be several reflection cuts. At one such cut
in block $b$, write $Q_S^b(F)$ for the symmetric quadratic form of
$\nu_S$. Assume that for $b\notin S$ and $|S|\le t$ there is a
nonnegative quadratic expression $D_S^b(F)$ such that
\begin{equation}
 Q_S^b(F)\ge-C_-\rho D_S^b(F),\qquad
 Q_{S\cup\{b\}}^b(F)\ge c_+D_S^b(F),
 \label{eq:abstractdom}
\end{equation}
where $C_-,c_+>0$ and $\rho>0$. Suppose terminal measures
$|S|=t+1$ are nonnegative at every cut in an unmarked block.

\begin{theorem}[Marking with a matched weight]
\label{thm:marking}
Choose $A\ge\max\{1,C_-/c_+\}$, put $a=A\rho$, and suppose $a\le1$.
Then
\begin{equation}
 M=\sum_{|S|\le t+1}a^{|S|}\nu_S
 \label{eq:abstractmixture}
\end{equation}
is reflection positive at every prescribed cut. If also
$\nu_\varnothing(1)\ge z_0>0$ and $\nu_S(1)\le B$, then
\[
 z_0\le Z:=M(1)\le2^\ell B.
\]
The total weight of nonleading components in either the normalized
full mixture or the normalized mixture with terminal components removed
is at most $2^\ell B a/z_0$. The normalized terminal mass satisfies
\begin{equation}
 \eta\le\binom\ell{t+1}\frac B{z_0}a^{t+1}.
 \label{eq:abstracteta}
\end{equation}
If every terminal mass is at least one, then also
\begin{equation}
 \eta\ge\binom\ell{t+1}\frac1{2^\ell B}a^{t+1}.
 \label{eq:abstractetalower}
\end{equation}
\end{theorem}
\begin{proof}
Fix a cut in block $b$. Pair each mark set not containing $b$ and of
size at most $t$ with its union with $\{b\}$. The exact decomposition is
\begin{align*}
 M={}&\sum_{\substack{b\notin S\\|S|\le t}}
 a^{|S|}(\nu_S+a\nu_{S\cup\{b\}})
 +a^{t+1}\sum_{\substack{b\notin S\\|S|=t+1}}\nu_S.
\end{align*}
Every term in a pair has quadratic form at least
$(ac_+-C_-\rho)D_S^b(F)\ge0$; the unpaired terms are terminal and
nonnegative. The mass statements follow by summing nonnegative masses.
In both retained mixtures the denominator is at least $z_0$, and
$a^{|S|}\le a$ for every nonempty retained $S$.
\end{proof}

Suppose all normalized component one-label marginals are the same
uniform law, and the leading conditional operator on mean-zero
one-label functions has norm at most $C_0\rho$. Revealing the component
index, applying conditional Jensen, and then using conditional-expectation
contraction in all other components gives
\begin{equation}
 \norm{T}_{L^2_0\to L^2}^2
 \le C_0^2\rho^2+\frac{2^\ell B}{z_0}a.
 \label{eq:abstractconditional}
\end{equation}
Here $T$ conditions one label on all the other labels under either
retained normalized mixture. Thus matching $a$ to a spectral error of
order $q^{-1/2}$ produces terminal mass $O(q^{-(t+1)/2})$ and
conditional norm $O(q^{-1/4})$. These exponents will be attained,
not only bounded, by the concrete construction.

\section{The finite permutation-block construction}
\label{sec:construction}

Fix integers $h\ge3$, $t\ge1$, and subgroup index $s\ge2$. Throughout
this section $q$ ranges over field orders for which $s\mid q-1$, and
$H_q\le\F_q^\times$ is the subgroup of index $s$. All constants may
depend on $h,t,s$, but not on $q$. The group $H_q$ is unique because
$\F_q^\times$ is cyclic.

\subsection{Vertices, cuts, and addresses}

Let $\Pi$ be the set of all permutation words of $[h]$. Its elements
are bijections from positions to symbols. Put
\begin{equation}
 K=h!,\qquad R=K/2,\qquad \ell=2t+1,\qquad
 V=\Pi^\ell,\qquad D=|V|=K^\ell,\qquad n=D/2.
 \label{eq:parameters}
\end{equation}
Fix an $h$-cycle $c$ on the positions. Right composition $w\mapsto wc$
shifts a word. For $0\le j<h$, let $v_j\in V$ have word $c^j$ in
every block, and set $v_h=v_0$.

A cut is specified by a block $b$ and distinct symbols $\alpha,\beta$.
Let $\theta$ interchange these symbols in block $b$, by composition on
the left. Define
\[
 V_+=\{v:\alpha\text{ occurs before }\beta\text{ in block }b\},
 \qquad V_-=\theta V_+.
\]
The involution has no fixed vertex and $|V_+|=|V_-|=n$.
A measure $\nu$ on $\F_q^V$ is reflection positive if, for every cut,
its matrix with row labels $z_{V_+}$ and column labels
$(z_{\theta v})_{v\in V_+}$ is symmetric positive semidefinite.
Equivalently its quadratic form
\[
 Q_\nu^b(F)=\int F(z_{V_+})F((z_{\theta v})_{v\in V_+})\,d\nu(z)
\]
is nonnegative for every real $F$, and the matrix is symmetric.

For a mark set $S\subseteq[\ell]$ of size $j\le t+1$, let
\[
 U_S=\Pi^{[\ell]\setminus S},\qquad m_j=R^j,
 \qquad u_S:V\to U_S
\]
be the projection that forgets marked blocks. At each address $u\in U_S$
choose an independent list $(Y_{u,i})_{1\le i\le m_j}$ of independent
uniform elements of $\F_q$. The total number of list entries is
\[
 |U_S|m_j=K^{\ell-j}R^j=D/2^j.
\]

For $j\le t$, draw a directed edge $u\to u'$ between addresses if
\begin{equation}
 \#\{a\notin S:u'_a=u_ac\}\ge t+1.
 \label{eq:addressedge}
\end{equation}
For each address edge require every pair of its list entries to satisfy
\begin{equation}
 Y_{u',i'}-Y_{u,i}\in H_q.
 \label{eq:listconstraints}
\end{equation}
Let $I_S(Y)$ be the indicator of all these requirements. At the terminal
level $j=t+1$, impose no field constraints, so $I_S=1$.

These list-constraint graphs have no loops: a permutation word cannot
be fixed by right multiplication by $c$. They have no opposite edges:
the blocks counted for opposite directions would be disjoint, since
a block counted for both would imply $c^2=1$. But two sets of $t+1$
blocks cannot fit in $\ell=2t+1$ blocks. Thus Theorem~\ref{thm:residuelists}
and its scalar version apply to every component and every induced
subgraph. Let $e_S$ be the number of edges of this fixed list graph.
For nonterminal $S$,
\begin{equation}
 \E_Y I_S(Y)=s^{-e_S}+O(q^{-1/2}).
 \label{eq:validitymass}
\end{equation}
In particular every such mass is positive for all sufficiently large
admissible $q$.

\subsection{A positive interaction among index choices}

Set $\delta=1/(10h^2)$. For distinct vertices $v,w\in V$ let
$J^S_{vw}=0$ unless $v,w$ differ by a symbol transposition in precisely
one marked block and agree in all other blocks. In that case, if the
two swapped positions are distance $g$ apart, set $J^S_{vw}=\delta^g$.
This is a symmetric nonnegative interaction.

Independently of the field lists, choose each index
$i_v\in[m_j]$ uniformly and independently. Give the index configuration
the unnormalized weight
\begin{equation}
 W_S(i)=\exp\left(\sum_{\{v,w\}\subset V}
                  J^S_{vw}1_{i_v=i_w}\right).
 \label{eq:indexweight}
\end{equation}
The sum is over unordered pairs. Define the nonnegative finite measure
$\nu_{q,S}$ by
\begin{equation}
 \int F(z)\,d\nu_{q,S}(z)
 =\E_Y\E_i I_S(Y)W_S(i)
           F\bigl((Y_{u_S(v),i_v})_{v\in V}\bigr).
 \label{eq:component}
\end{equation}
This definition is a pushforward of a weighted finite product measure;
it does not presume that different index configurations give different
label configurations.

Each vertex has at most $j\binom h2$ interaction neighbors, each of
weight at most $\delta$. Hence
\begin{equation}
 1\le W_S\le B,
 \qquad B=\exp\left(\frac{D(t+1)\binom h2\delta}{2}\right),
 \qquad 0\le\nu_{q,S}(1)\le B.
 \label{eq:B}
\end{equation}
Terminal components have mass between $1$ and $B$, independent of
$q$. In the leading component $S=\varnothing$, every list has length
one and $W_\varnothing=1$. Its density relative to independent uniform
labels at all vertices is the indicator $I_\varnothing$.

Every component is invariant under every reflection $\theta$.
Indeed left symbol permutations preserve positional gaps and commute
with the right position shift $c$. If the cut block is unmarked,
$\theta$ permutes addresses and therefore whole independent lists;
if it is marked, the address of $v$ equals that of $\theta v$.
In both cases the index relabeling $i_v\mapsto i_{\theta v}$ preserves
the product law and its weight. This proves symmetry of all cut
matrices. It does not by itself prove positivity.

\subsection{Strict positivity at a marked cut}

\begin{lemma}[Marked-cut domination]
\label{lem:marked}
Let $b\in S$, $|S|=j$, and use a cut in block $b$. Put
\[
 \kappa=\frac9{10}\delta^{h-1},\qquad
 \lambda=(e^\kappa-1)^n,\qquad
 c_* =\lambda R^{-(t+1)n}.
\]
Then every real half-label function satisfies
\begin{equation}
 Q_{q,S}^b(F)
 \ge\lambda m_j^{-n}\E_{Y,i_+}I_S(Y)
           F((Y_{u_S(v),i_v})_{v\in V_+})^2.
 \label{eq:marked}
\end{equation}
In particular the coefficient is at least $c_*>0$.
\end{lemma}
\begin{proof}
Index a real symmetric matrix $A$ by $x,y\in V_+$ and set
$A_{xy}=J^S_{x,\theta y}$. Suppose that the cut symbols in $x_b$
occupy positions $i<j$ with gap $g=j-i$. Then $A_{xx}=\delta^g$.
An off-diagonal transposition taking $x$ across this cut either moves
the first cut symbol beyond position $j$ or the second before position
$i$. Its gap is at least $g+1$. No transposition in another block can
cross the cut. There are fewer than $h^2$ possible transpositions, so
\[
 \sum_{y\ne x}|A_{xy}|\le h^2\delta^{g+1}
                         =\tfrac1{10}A_{xx}.
\]
Using $2|a_xa_y|\le a_x^2+a_y^2$ gives
$A\succeq\kappa I_n$.

Fix the field lists, and write $m=m_j$ and
$\mathcal I=[m]^{V_+}$. For an index configuration $i\in\mathcal I$,
let $u(i)$ be its one-hot vector in $\R^{V_+\times[m]}$. The cross
interaction matrix between the two half index configurations is
\[
 C(i,k)=\exp\bigl(u(i)^T(A\otimes I_m)u(k)\bigr).
\]
There is no factor two: a cross unordered pair has the unique form
$\{x,\theta y\}$ with $x,y\in V_+$.
Decompose this entrywise exponential into the product of the factors
from $\kappa I_n$ and $A-\kappa I_n$. The latter factor, denoted by
$P$, is positive semidefinite. To see this, express
$u(i)^T((A-\kappa I_n)\otimes I_m)u(k)$ as an inner product; every
term in its exponential series is a Gram matrix of tensor powers.
Also $P(i,i)\ge1$.

The first factor is the matrix
\[
 E=\bigotimes_{x\in V_+}
       \bigl((e^\kappa-1)I_m+\one\one^T\bigr)
       \succeq\lambda I_{\mathcal I}.
\]
The entrywise product of two Gram matrices is again a Gram matrix,
using tensor products of their representing vectors. Therefore
\[
 C=(E-\lambda I)\circ P+\lambda(I\circ P)\succeq\lambda I.
\]
Internal plus-side index interactions multiply this matrix on the
left and right by the same diagonal matrix with diagonal entries at
least one. The lower bound $\lambda I$ is preserved.

Because $b$ is marked, reflected vertices have the same address.
The half-label function thus produces the same vector $f_Y(i)$ on
both sides, even if some emitted field labels coincide. The uniform
double index average is at least
\[
 \lambda m^{-2n}\sum_{i\in\mathcal I} f_Y(i)^2
 =\lambda m^{-n}\E_{i_+}f_Y(i_+)^2.
\]
Multiply by the nonnegative $I_S(Y)$ and average over $Y$.
Finally $m_j\le R^{t+1}$ gives the uniform coefficient $c_*$.
\end{proof}

\section{Matching the spectral error and normalizing the laws}
\label{sec:normalization}

\subsection{The improved error at an unmarked cut}

If $b\notin S$, the addresses split into two independent halves.
Let $I_+(Y_+)$ impose the field constraints internal to plus addresses,
and put
\[
 D_{q,S}^b(F)=\E_{Y_+,i_+}I_+(Y_+)
            F((Y_{u_S(v),i_v})_{v\in V_+})^2.
\]

\begin{lemma}[Unmarked error and next-level comparison]
\label{lem:unmarked}
For $b\notin S$, $|S|\le t$, one has
\begin{align}
 Q_{q,S}^b(F)&\ge-2B^2q^{-1/2}D_{q,S}^b(F),
       \label{eq:unmarkednew}\\
 Q_{q,S\cup\{b\}}^b(F)&\ge c_*R^{-n}D_{q,S}^b(F).
       \label{eq:nextlevel}
\end{align}
For $b\notin S$, $|S|=t+1$, the form $Q_{q,S}^b$ is nonnegative.
\end{lemma}
\begin{proof}
For the first assertion, the field indicator factors as
\[
 I_S(Y_+,Y_-)=I_+(Y_+)I_+(Y_-)K_b(Y_+,Y_-),
\]
after using reflection to identify the minus list space with a copy of
the plus list space. The crossing kernel is a simple bipartite residue
pattern. By \eqref{eq:subgrouplists}, if it has $e_b$ edges,
\[
 \norm{K_b-s^{-e_b}}_{\op}\le2q^{-1/2}.
\]
All index interactions change marked blocks and hence preserve the
unmarked cut block. Consequently the index weight splits into identical
plus and reflected-minus factors. Define
\[
 G(Y_+)=I_+(Y_+)\E_{i_+}W_+(i_+)
                   F((Y_{u_S(v),i_v})_{v\in V_+}).
\]
Then $Q_{q,S}^b(F)=\langle G,T_{K_b}G\rangle$. The constant model
contributes $s^{-e_b}(\E G)^2\ge0$, and
$\norm G_2^2\le B^2D_{q,S}^b(F)$ by Cauchy--Schwarz in the index
average. This proves \eqref{eq:unmarkednew}. If $S$ is terminal,
$K_b=1$, so the same expression is a nonnegative square.

For the second assertion, write $S'=S\cup\{b\}$ and $|S|=j$.
There are $R$ words in the cut's plus half of $\Pi$. Group the old
plus-side lists above one new address into a single list of length
$Rm_j=m_{j+1}$. This is an identification of independent uniform
field-entry collections, not a conditioning event. Under it, old
internal validity implies new validity: any new address edge has
$t+1$ matching clean blocks outside $b$, so every old plus-address
pair above it was already constrained. For a terminal new component
the implication is automatic.

Apply Lemma~\ref{lem:marked} to $S'$. Restrict each new half index to
the sublist corresponding to that vertex's actual plus word in block
$b$. The restriction has probability exactly $R^{-n}$ because the
$n$ indices are independent. Conditional on it, the remaining indices
are independent uniform old indices and give the old half labels.
The integrand is a square, so restriction only decreases it. The
validity implication then gives \eqref{eq:nextlevel}.
\end{proof}

\subsection{The optimized mixture}

Define the fixed positive constants
\begin{equation}
 c_+=c_*R^{-n},\qquad
 A=\max\{1,2B^2/c_+\},\qquad a_q=Aq^{-1/2}.
 \label{eq:optimizedweight}
\end{equation}
Only field orders large enough that $a_q\le1$ are considered from now
on. Set
\begin{align}
 M_q&=\sum_{|S|\le t+1}a_q^{|S|}\nu_{q,S},& Z_q&=M_q(1),
       \label{eq:concretemixture}\\
 \mu_q&=M_q/Z_q,&
 \varepsilon_q&=Z_q^{-1}a_q^{t+1}\sum_{|S|=t+1}\nu_{q,S},
       \label{eq:mu-epsilon}\\
 \eta_q&=\varepsilon_q(1),&
 \mu_q^\circ&=(\mu_q-\varepsilon_q)/(1-\eta_q).
       \label{eq:nonexceptional}
\end{align}
The leading mass tends to $s^{-e_0}$, where $e_0=e_\varnothing$,
by \eqref{eq:validitymass}. Thus for sufficiently large $q$,
\begin{equation}
 \nu_{q,\varnothing}(1)\ge z_0:=\tfrac12s^{-e_0}>0,
 \qquad z_0\le Z_q\le2^\ell B.
 \label{eq:z0}
\end{equation}
Lemmas~\ref{lem:unmarked} and \ref{lem:marked} verify every premise
of Theorem~\ref{thm:marking}, with $\rho=q^{-1/2}$.
It follows that $\mu_q$ is reflection positive at every prescribed cut.
The submeasure $\varepsilon_q$ is nonnegative, and $\eta_q\to0$;
therefore both normalized laws are well defined.

\begin{proposition}[Exact order of the exceptional mass]
\label{prop:eta}
For terminal $S$ put
$\beta_S=\E_i W_S(i)$ and $\beta=\sum_{|S|=t+1}\beta_S$.
These are finite positive constants independent of $q$. Then
\begin{equation}
 \eta_q=A^{t+1}\beta s^{e_0}q^{-(t+1)/2}
                       \bigl(1+O(q^{-1/2})\bigr).
 \label{eq:etaasymptotic}
\end{equation}
In particular $\eta_q=\Theta(q^{-(t+1)/2})$.
\end{proposition}
\begin{proof}
Terminal components have $I_S=1$, so their raw total mass is exactly
$\beta a_q^{t+1}$. The nonleading contribution to $Z_q$ is $O(a_q)$
by \eqref{eq:B}. The leading mass is
$s^{-e_0}+O(q^{-1/2})$. Hence
$Z_q=s^{-e_0}+O(q^{-1/2})$, whose reciprocal has the corresponding
expansion. Substitute into \eqref{eq:mu-epsilon}.
\end{proof}

\subsection{Uniform marginals and the exact selected-pair law}

Each component is invariant under simultaneous changes
$z_v\mapsto az_v+b$ with $a\in H_q$, $b\in\F_q$. Apply this affine
change to every field-list entry in \eqref{eq:component}; it preserves
the uniform product law, all difference constraints, and all index
weights. In particular every positive-mass normalized component has
uniform single-label marginals, because translations act transitively.
So do both laws in \eqref{eq:mu-epsilon}--\eqref{eq:nonexceptional}.

For nonterminal $S$, at least $\ell-t=t+1$ clean blocks remain.
Each selected consecutive pair $v_j,v_{j+1}$ is related by $c$ in every
one of these blocks. Its addresses form an edge, and every pair of
entries in their lists is constrained. Therefore
\[
 z_{v_{j+1}}-z_{v_j}\in H_q\quad(0\le j<h)
\]
holds throughout the support of $\mu_q^\circ$.
Let
\[
 \Omega_q=\{(x,y)\in\F_q^2:y-x\in H_q\}.
\]
The affine group with multipliers in $H_q$ acts transitively on
$\Omega_q$: to carry $(x,y)$ to $(x',y')$, take
$a=(y'-x')/(y-x)\in H_q$ and $b=x'-ax$. Invariance and support
therefore prove that the selected pair is exactly uniform on $\Omega_q$
under $\mu_q^\circ$.

At the terminal level, $t\ge1$ clean blocks still remain. Consecutive
selected vertices have distinct addresses, so their output labels lie
in independent uniform lists. There are no field constraints, and
the index weight is independent of the field entries. Their pair law
is consequently uniform on $\F_q^2$, even after averaging all indices.
For the full law the selected pair distribution is exactly
\begin{equation}
 (1-\eta_q)\Unif(\Omega_q)+\eta_q\Unif(\F_q^2).
 \label{eq:pairlaw}
\end{equation}
No reflection-positivity assertion is made for $\mu_q^\circ$; removing
a positive submeasure need not preserve positive semidefiniteness.

\section{Sharp conditional prediction estimates}
\label{sec:mixing}

For a law $\rho$ equal to either $\mu_q$ or $\mu_q^\circ$, and a
vertex $v$, define
\[
 T_{\rho,v}f
 =\E_\rho\bigl[f(z_v)\mid(z_w)_{w\ne v}\bigr].
\]
The input space $L^2_0(\F_q)$ is the space of uniform-mean-zero
complex functions. The output norm uses the actual other-label
marginal of $\rho$.

\subsection{An upper bound without uniform conditional densities}

\begin{lemma}[Leading-component prediction]
\label{lem:leadingconditional}
For the normalized leading component $\rho_0$, uniformly over $v$,
\[
 \norm{T_{\rho_0,v}}_{L^2_0\to L^2}=O(q^{-1/2}).
\]
\end{lemma}
\begin{proof}
Write $x=z_v$, $y=(z_w)_{w\ne v}$. The leading density is
$L(y)K(x,y)$, where $K$ is the product of the $m_v$ constraints incident
to $v$ and $L$ is the indicator of all other constraints. Put
\[
 c_v=s^{-m_v},\quad a(y)=\E_xK(x,y),\quad
 b_f(y)=\E_x f(x)K(x,y).
\]
If $m_v=0$, $x$ is independent and the conclusion is immediate.
Otherwise Theorem~\ref{thm:residuelists}, using transposition as needed,
gives
\begin{equation}
 \norm{a-c_v}_2\le\Delta,\qquad
 \norm{b_f}_2\le\Delta\norm f_2,
 \qquad \Delta=2q^{-1/2}.
 \label{eq:ab}
\end{equation}
Let $Z_0=\nu_{q,\varnothing}(1)\ge z_0$. The actual other-label
marginal has density $L(y)a(y)/Z_0$. Thus
\begin{equation}
 \norm{T_{\rho_0,v}f}_2^2
 =Z_0^{-1}\E_y L(y)\frac{|b_f(y)|^2}{a(y)},
 \label{eq:conditionalratio}
\end{equation}
with zero assigned where $a=0$. On $a\ge c_v/2$, the expression is
bounded by $2z_0^{-1}c_v^{-1}\Delta^2\norm f_2^2$ after integration.
On its complement, Chebyshev gives probability at most
$4\Delta^2/c_v^2$. Weighted Cauchy--Schwarz gives the pointwise estimate
\[
 \frac{|b_f(y)|^2}{a(y)}
 \le\E_x|f(x)|^2K(x,y)\le\norm f_2^2.
\]
Since $0\le L\le1$, the complement contributes at most
$4z_0^{-1}c_v^{-2}\Delta^2\norm f_2^2$.
All $c_v$ are fixed positive numbers, and there are finitely many
vertices. Combining the estimates proves the result uniformly in $v$.
\end{proof}

By \eqref{eq:abstractconditional}, with all normalized components
having uniform input marginals, Lemma~\ref{lem:leadingconditional}
implies
\begin{equation}
 \norm{T_{\rho,v}}_{L^2_0\to L^2}^2
 \le C_1q^{-1}+C_2a_q=O(q^{-1/2}).
 \label{eq:conditionalupper}
\end{equation}
Revealing the latent component index is legitimate even though the
index cannot in general be recovered from the output field labels.
Conditional Jensen only increases the squared prediction norm.

\subsection{A collision lower bound}

\begin{lemma}[Detectable triple collision]
\label{lem:collision}
Let $X$ be uniform on $\F_q$, and let the observation $Y$ include
two field-valued variables $W,U$. Put $E=\{W=U\}$.
Suppose, for some $a>0$ and fixed $c,C>0$,
\[
 \Prob(E)\le Ca,\qquad \Prob(E\cap\{X=W\})\ge ca.
\]
Then, for all sufficiently large $q$ depending only on $c,C$,
\begin{equation}
 \norm{f\mapsto\E[f(X)\mid Y]}_{L^2_0(\F_q)\to L^2}
 \ge\frac{c}{2\sqrt C}\sqrt a.
 \label{eq:collisionlower}
\end{equation}
\end{lemma}
\begin{proof}
For each nontrivial additive character $\psi$ use the mean-zero unit
input $f_\psi(x)=\psi(x)$ and the observation test
$g_\psi(Y)=\psi(W)1_E$. Its squared norm is $\Prob(E)$.
Orthogonality of all $q$ additive characters gives
\[
 \frac1{q-1}\sum_{\psi\ne1}\psi(X-W)
 =\frac q{q-1}1_{X=W}-\frac1{q-1}.
\]
The average real part of
$\E[f_\psi(X)\overline{g_\psi(Y)}]$ is therefore at least
$(qca-Ca)/(q-1)$, which is at least $ca/2$ for sufficiently large
$q$. Some character attains at least this real part. Pairing its
conditional expectation with $g_\psi$ and dividing by
$\norm{g_\psi}_2\le\sqrt{Ca}$ proves \eqref{eq:collisionlower}.
\end{proof}

\begin{theorem}[Sharp prediction rate for the optimized hierarchy]
\label{thm:sharpconditional}
For the construction above, every vertex $v$ and both
$\rho\in\{\mu_q,\mu_q^\circ\}$ satisfy
\begin{equation}
 \boxed{\qquad
 \norm{T_{\rho,v}}_{L^2_0\to L^2}=\Theta(q^{-1/4}).
 \qquad}
 \label{eq:sharpconditional}
\end{equation}
The positive constants in the two bounds can be chosen uniformly
over all vertices and both laws.
\end{theorem}
\begin{proof}
The upper bound is \eqref{eq:conditionalupper}. For the lower bound,
fix $v$ and a block $b$. Choose distinct vertices $w,u$, both different
from $v$, which agree with $v$ outside $b$. Such choices exist because
$K=h!\ge6$. Use the observation event $E=\{z_w=z_u\}$.

The leading unnormalized measure has density bounded by one relative
to independent uniform field labels, so its mass on $E$ is at most
$1/q$. All nonleading components together have normalized mass
$O(a_q)$, in either law. Their normalizers are at least $z_0$, so
\begin{equation}
 \Prob_\rho(E)\le C(q^{-1}+a_q)\le C'a_q.
 \label{eq:collisionupperE}
\end{equation}

In the component $S=\{b\}$ the vertices $v,w,u$ have the same
address and share a list of length $R$. The independent raw index
event $i_v=i_w=i_u$ has probability $R^{-2}$. On it all three output
labels agree. It is independent of the field-list validity event,
and $W_S\ge1$. Consequently
\[
 \nu_{q,\{b\}}(E\cap\{z_v=z_w\})
 \ge R^{-2}\E_YI_{\{b\}}(Y)\ge c_1>0
\]
for sufficiently large $q$, by \eqref{eq:validitymass}. The constant
can be uniform over the finitely many blocks. Since $t\ge1$, this
singleton mark set is nonterminal and occurs in both laws with raw
coefficient $a_q$. Both normalizers are at most $2^\ell B$. Thus
\begin{equation}
 \Prob_\rho(E\cap\{z_v=z_w\})\ge c_2a_q.
 \label{eq:collisionjoint}
\end{equation}
Apply Lemma~\ref{lem:collision} with $X=z_v$, $W=z_w$, $U=z_u$.
The input is uniform and $E$ is observable without $z_v$. The result
is $\norm{T_{\rho,v}}\ge c_3\sqrt{a_q}=c_3\sqrt A\,q^{-1/4}$.
\end{proof}

The lower bound identifies a specific obstruction. A singleton marked
component has probability of order $q^{-1/2}$, and it creates triple
collisions which an observer can partially detect from the other two
labels. Concentrating the observation test on these collisions costs
the square root of that probability. Increasing the terminal marking
depth does not remove singleton components, so it cannot improve this
prediction exponent within the same hierarchy.

\Needspace{12\baselineskip}
\section{The 17-block theorem and its quantitative meaning}
\label{sec:17blocks}

\begin{theorem}[Compressed square-difference tuple laws]
\label{thm:compressed}
Set $h=24$, $t=8$, $\ell=17$, and $V=\Pi^{17}$, with selected vertices
and reflection cuts as in Section~\ref{sec:construction}. There is a
fixed integer $P$ such that every prime $p\ge P$ admits a probability
law $\mu_p$ on $\F_p^V$ and a nonnegative submeasure
$\varepsilon_p\le\mu_p$ with mass $\eta_p\le p^{-4}$ for which:
\begin{enumerate}
\item $\mu_p$ is reflection positive at every specified cut.
\item Both $\mu_p$ and the normalized nonexceptional law $\mu_p^\circ$
      have uniform individual marginals and are invariant under common
      translations and nonzero-square multiplications of all labels.
      The submeasure $\varepsilon_p$ has the same affine invariance.
\item On the support of $\mu_p^\circ$ the designated directed
      $24$-cycle has nonzero-square differences. Each selected ordered
      pair is uniform on $\Omega_p$ under $\mu_p^\circ$, and has
      the exact mixture law \eqref{eq:pairlaw} under $\mu_p$.
\item $\eta_p=C_{\rm exc}p^{-9/2}(1+O(p^{-1/2}))$ for the explicit
      finite positive constant $C_{\rm exc}=A^9\beta2^{e_0}$.
\item For each $v$ and either law,
      $\norm{T_{\rho,v}}=\Theta(p^{-1/4})$. In particular, for every
      fixed $0<\gamma<1/4$, increasing $P$ if necessary gives
      $\norm{T_{\rho,v}}\le p^{-\gamma}$ simultaneously for all
      vertices and both laws.
\end{enumerate}
\end{theorem}
\begin{proof}
Use the construction with subgroup index $s=2$. Positivity and all
marginal assertions were established in Section~\ref{sec:normalization}.
Proposition~\ref{prop:eta} gives the exceptional asymptotic, hence
$\eta_p\le p^{-4}$ for large $p$. Theorem~\ref{thm:sharpconditional}
gives the last assertion. If $\norm{T_{\rho,v}}\le C p^{-1/4}$,
then $p\ge C^{1/(1/4-\gamma)}$ suffices for the constant-one bound.
All relevant finite constants can be maximized to choose one $P$.
\end{proof}

\begin{table}[htbp]
\centering\small
\begin{tabularx}{\textwidth}{@{}lXX@{}}
\toprule
Quantity & Source construction \cite{square} & Construction proved here\\
\midrule
Cross-list error & $(32m)^{1/4}p^{-1/8}$ & $2(1-2^{-m})\Delta_p$\\
Marking weight per block & $p^{-1/16}$ & $A p^{-1/2}$, fixed $A$\\
Last nonterminal depth & $80$ & $8$\\
Number of blocks & $161$ & $17$\\
Number of labels & $(24!)^{161}$ & $(24!)^{17}$\\
Terminal mass estimate & $O(p^{-81/16})$, used to give $p^{-4}$ &
  $\Theta(p^{-9/2})$, still at most $p^{-4}$\\
Prediction estimate & $\norm T\le p^{-1/64}$ beyond a threshold &
  $\norm T=\Theta(p^{-1/4})$; every constant-one exponent $\gamma<1/4$\\
Selected pair law & Exact square-pair law after terminal removal & Same property\\
\bottomrule
\end{tabularx}
\caption{A comparison of proved local parameters. The smaller construction
retains the required terminal-mass target; it does not retain the source's
stronger unused exponent $81/16$. Constants and starting field orders change.}
\label{tab:comparison}
\end{table}

The ratio of the two label counts is exactly $(24!)^{144}$.
This is a substantial reduction in a finite construction, but both
counts and the resulting constants are still enormous. In particular
$A$ contains the reciprocal of a very small marked-cut eigenvalue;
the theorem is not a practical sampler for a realistically sized tuple.

\begin{corollary}[Depth versus terminal error]
\label{cor:tradeoff}
For every fixed $h\ge3$, fixed subgroup index $s\ge2$, and integer
$t\ge1$, the construction on $(h!)^{2t+1}$ labels has
\[
 \eta_q=\Theta(q^{-(t+1)/2}),\qquad
 \norm{T_{\rho,v}}=\Theta(q^{-1/4}).
\]
For any desired terminal bound $q^{-B_0}$, it suffices to choose
$t$ with $(t+1)/2>B_0$. For $B_0=4$, $t=8$ gives 17 blocks.
\end{corollary}

The strict inequality in this sufficient condition absorbs the finite
leading constant in \eqref{eq:etaasymptotic}. We do not claim that
17 is the least possible number of blocks among all constructions.
Nor do we rule out a depth with $(t+1)/2=B_0$ if its leading constant
can be made small enough. What is sharp here is the exponent for the
specified fixed positive-weight hierarchy.

\subsection{What this changes in the integer problem}

The full square-difference argument in \cite{square} proceeds beyond
the tuple law through additional multilinear and arithmetic estimates.
Theorems~\ref{thm:residuelists} and \ref{thm:compressed} replace two
finite ingredients with stronger ones. The parameter $h=24$ is retained
to preserve the designated cycle used in that argument; the tuple
theorem alone permits every $h\ge3$.

An improvement in a local exponent does not automatically identify
an improved value of $c$ in $|A|\le C N^{1-c}$. Deriving such a value
would require carrying the new degree, thresholds, conditional exponent,
and all other losses through the remaining proof. This report establishes
the precise local replacements and their limits, and leaves that global
optimization as a separate research problem. Its finite theorems remain
valid independently of that global analysis.

\section{Reproducible verification and a formalization route}
\label{sec:verification}

The proofs above are analytic and finite-dimensional; no numerical
observation is used to establish a theorem. A companion program,
\texttt{code/spectral\_checks.py}, tests normalization, directions,
edge ordering, and small marked matrices. Its recorded output is
\texttt{data/spectral\_validation.json}. The deterministic seed is
\texttt{20261007159}.

\begin{table}[htbp]
\centering\small
\begin{tabularx}{\textwidth}{@{}lXr@{}}
\toprule
Type & Check & Count\\
\midrule
Exact rational & Telescoping identities at all selected coordinate assignments & 1,024\\
Exact rational & Optimal-order comparisons against permutations & 5,913\\
Exact rational & Marked cross-interaction diagonal dominance, $h=3,4,5$ & 3 matrices\\
Floating-point & Square-kernel Fourier and singular-value norms, primes $\le101$ & 25 primes\\
Floating-point & Every oriented $2$-by-$2$ pattern at $p=3,5,7$ & 243 patterns\\
Floating-point & Heterogeneous kernels with nonuniform base probabilities & 28 cases\\
Floating-point & Permutation costs in those heterogeneous cases & 23,652 orders\\
\bottomrule
\end{tabularx}
\caption{Supporting verification. Exact checks use rational arithmetic;
the floating-point checks are diagnostics, with absolute tolerance
$2\cdot10^{-10}$, and are not proof certificates for spectral inequalities.}
\end{table}

The maximum observed absolute discrepancy in the exact square-norm
formula across the numerical tests was below $9\cdot10^{-17}$.
The same program includes the repeated-edge obstruction, which is
supposed to violate the conclusion after its simplicity hypothesis is
removed. The huge $24$-symbol tuple measure is not enumerated.
Its positivity is proved by the uniform matrix argument in
Lemma~\ref{lem:marked}, and its prediction lower bound is proved
without constructing its full probability table.

A useful Lean development would start with finite weighted probability
spaces, an explicit bilinear formulation of the spectral norm, and
Lemma~\ref{lem:exposed}. The telescope and ordering rule then become
finite sums and products. A second module can supply additive-character
orthogonality and the Gauss calculation. The marked construction should
be handled at the level of pushforward finite measures and Gram matrices,
so that repeated field labels are never treated as distinct states by
assumption. The final module would formalize the latent-mixture Jensen
bound and Lemma~\ref{lem:collision}. No Lean files in this package claim
these steps have been completed.

\section{Further research questions}
\label{sec:questions}

\begin{question}[Pattern-specific spectral constants]
For a fixed simple bipartite pattern $H$, determine the smallest
$C_H(d)$ for which $\norm{K_W-d^{|E(H)|}}_{\op}\le C_H(d)\epsilon$
holds for all one-bounded kernels with edge errors at most $\epsilon$.
The geometric bound is uniform in $H$, but need not be optimal for a
particular matching, tree, or highly symmetric pattern.
\end{question}
The one-edge example settles the best possible universal power of the
error. It does not settle pattern-dependent constants or equality cases.
An additional exact-degree assumption, $W_e\one=d\one$ and
$W_e^*\one=d\one$, may improve the constants by making different
first-order edge contributions orthogonal.

\begin{question}[Smaller positivity constants]
Can the marked-cut eigenvalue $\lambda R^{-(t+2)n}$ be replaced by a
much larger explicit lower bound on the subspace actually reached by
field-label functions? The current proof bounds the entire index matrix,
including directions which may never be observed after pushforward.
\end{question}
This could lower $A$ and the starting field order substantially while
keeping the same exponents. The question concerns a finite matrix
compression, and is suitable for symbolic representation-theoretic study.

\begin{question}[Avoiding the collision obstruction]
Is there a different reflection-positive coupling with the same exact
selected-pair law and terminal error $O(q^{-4})$, but conditional norm
$O(q^{-1/4-\varepsilon})$ for some $\varepsilon>0$?
\end{question}
Theorem~\ref{thm:sharpconditional} rules out obtaining this merely by
sharpening the upper-bound analysis of the present fixed mixture.
A successful construction must suppress or hide the observable triple
collisions used in Lemma~\ref{lem:collision}, or avoid singleton marked
mass of order $q^{-1/2}$.

\begin{question}[A minimal block architecture]
How few permutation blocks can preserve all the required reflections,
the directed cycle, and an exceptional mass at most $q^{-4}$?
Can the geometric pairing of mark sets be replaced by a nonuniform
family with fewer levels or a different incidence design?
\end{question}
The 17-block construction is a proved sufficient design, not a universal
lower bound. Different positive weights for different mark sets may
lead to a finite optimization problem constrained by all reflection cuts.

\begin{question}[Complete arithmetic propagation]
After substituting Theorem~\ref{thm:compressed} into the remaining
square-difference argument, what explicit global exponent and threshold
can be proved? Which remaining step dominates the quantitative loss?
\end{question}
This requires an actual derivation through all subsequent inequalities.
The local degree and conditional improvements should be propagated
together; improving one and leaving the other at its old value could
hide the available gain.

\begin{question}[Growing subgroup index and pattern size]
Can the fixed-$s$, fixed-pattern tuple theorem be made uniform when the
subgroup index or the number of constraints grows with $q$?
\end{question}
Equation~\eqref{eq:subgrouplists} remains a uniform absolute-error
estimate. The obstacle is the diminishing normalizing mass and the
growth of finite interaction constants. A useful uniform theorem must
control these explicitly, rather than absorb them into a fixed threshold.

\begin{question}[More general constraint spaces]
Which finite homogeneous relations have both a small centered spectral
norm and a transitive symmetry group on their allowed ordered pairs,
so that the same construction gives exact pair marginals?
\end{question}
The replacement theorem itself applies to arbitrary finite probability
spaces. Extending the tuple theorem requires the additional symmetry
and positivity-of-counts inputs. Affine finite-field relations and
other association schemes provide concrete candidates.

\appendix
\section{Source and dependency ledger}

The exact source locations used in the comparison are listed below.
Paths in the OpenAI repository are relative to
\nolinkurl{preprints/A-power-saving-for-square-difference-free-sets-September-24-2026/}.

\begin{table}[htbp]
\centering\small
\begin{tabularx}{\textwidth}{@{}lX@{}}
\toprule
Source location & Use in this report\\
\midrule
\texttt{README.md} & Manuscript title, date, and author attribution.\\
\texttt{build/sections/graph.tex} & Original single-square calculation and
  the list-mixing estimate \eqref{eq:old-bound}. The new operator proof
  is Section~\ref{sec:replacement}.\\
\texttt{build/sections/tuples.tex} & Permutation blocks, index interaction,
  marked-cut positivity, and list regrouping. The construction is
  reproved with parameters $h,t,s$ in Sections~\ref{sec:construction}--\ref{sec:normalization}.\\
\texttt{lean/docs/182.md} (repository root) & The source's stated
  formalization scope; it is not a formalization of the new refinements.\\
ProveIt Reports 288--299 intake records & Context for sharp finite
  refinements and the distinction between proved local interfaces and
  global conclusions. No theorem from them is a premise of the present proof.\\
\bottomrule
\end{tabularx}
\end{table}

The source proof of the marked matrix estimate has been retained in
mathematical substance with attribution. The key changes are: retaining
arbitrary $L^2$ list tests in the exposed-edge step; retaining geometric
constant factors in the telescope; using a fixed constant times the
actual spectral scale as the marking weight; allowing arbitrary fixed
subgroup index; and proving the observable-collision lower bound.

For a concise dependency check, the graph theorem uses only finite
product measure, Cauchy--Schwarz, and a one-edge norm hypothesis.
The finite-field specialization adds elementary character orthogonality.
The tuple theorem adds finite Gram-matrix positivity, exact regrouping
of independent lists, and normalization by a positive counted mass.
The sharp lower bound adds only a three-index coincidence event and
the orthogonality identity for additive characters. No global
Szemer\'edi or square-difference theorem is imported into these proofs.

\begin{thebibliography}{99}
\raggedright\small
\bibitem{catalogue}
OpenAI, \emph{Mathematics manuscript collection}, October 2026 release.
\url{https://github.com/openai/math}. Snapshot \sourcecommit.

\bibitem{square}
OpenAI, \emph{A power saving for square-difference-free sets},
manuscript dated 24 September 2026, OpenAI Math Release, publicly
released 6 October 2026.
\url{https://github.com/openai/math/tree/adc7f1241b42e322a6451854ab7e4b4c146bf78a/preprints/A-power-saving-for-square-difference-free-sets-September-24-2026}.

\bibitem{proveit}
V. Reshetnikov and contributors, \emph{ProveIt}, repository snapshot
\texttt{fe7165a3cda1f9182b2a258cf13431a906c4b0c5}, 7 October 2026.
\url{https://github.com/VladimirReshetnikov/ProveIt}.

\bibitem{energyintake}
ProveIt contributors, \emph{Intake of Reports 288--293: hereditary
relative energy}, commit \texttt{487859db599dfd3a20d24b11595a44973047babf},
7 October 2026.
\url{https://github.com/VladimirReshetnikov/ProveIt/commit/487859db599dfd3a20d24b11595a44973047babf}.

\bibitem{refinementsintake}
ProveIt contributors, \emph{Intake of Reports 294--299 and product-property
refinements}, commit \texttt{bbde2b873b2dc7b48af1730926452bdda18bf9c8},
7 October 2026.
\url{https://github.com/VladimirReshetnikov/ProveIt/commit/bbde2b873b2dc7b48af1730926452bdda18bf9c8}.

\bibitem{zhao}
Y. Zhao, \emph{Graph Theory and Additive Combinatorics: Exploring
Structure and Randomness}, Cambridge University Press, 2023.
Author's manuscript and resources: \url{https://yufeizhao.com/gtacbook/}.

\bibitem{cfz}
D. Conlon, J. Fox, and Y. Zhao, \emph{Extremal results in sparse
pseudorandom graphs}, Advances in Mathematics \textbf{256} (2014),
206--290. \url{https://arxiv.org/abs/1204.6645}. See Section 3 for
edge-exposure counting arguments.

\bibitem{lss}
J. Leng, A. Sah, and M. Sawhney, \emph{Improved bounds for Szemer\'edi's
theorem}, 2024. \url{https://arxiv.org/abs/2402.17995}.

\bibitem{bloommaynard}
T. F. Bloom and J. Maynard, \emph{A new upper bound for sets with no
square differences}, Compositio Mathematica \textbf{158} (2022),
1777--1798. \url{https://arxiv.org/abs/2011.13266}.
\end{thebibliography}
\end{document}
